% Part: first-order-logic
% Chapter: completeness
% Section: compactness

\documentclass[../../../include/open-logic-section]{subfiles}

\begin{document}

\iftag{FOL}
      {\olfileid{fol}{com}{com}}
      {\olfileid{pl}{com}{com}}

\olsection{Teorema Kekompakan}

Salah sawijining akibat wigati saka teorema kelengkapan yaiku teorema
kekompakan. Teorema kekompakan nyatakake yen saben subhimpunan
\emph{winates} saka sawijining himpunan !!{sentence}s bisa dicukupi,
mula kabeh himpunan kasebut bisa dicukupi---sanajan himpunan kasebut
dhewe tanpa wates. Bab iki ora cetha babar pisan. Ing pandeleng
kapisan paling ora, katon ora ana apa-apa kang nyegah anane himpunan
!!{sentence}s tanpa wates kang kontradiktif, nanging kontradiksi mau,
upamane, mung tuwuh saka cacahe kang tanpa wates. Teorema kekompakan
nyatakake yen kahanan kaya mangkono bisa disingkirake: ora ana
himpunan !!{sentence}s tanpa wates kang ora bisa dicukupi nanging
saben subhimpunan winatese bisa dicukupi. Kaya teorema kelengkapan,
teorema iki duwe versi kang gegayutan karo akibat semantis: yen
sawijining himpunan !!{sentence}s tanpa wates ngakibatake sawijining
pratelan, sawijining subhimpunan winates wae wis ngakibatake pratelan
kasebut.

\begin{defn}
  Himpunan $\Gamma$ saka !!{formula}s diarani \emph{bisa dicukupi
  kanthi winates} yen lan mung yen saben $\Gamma_0 \subseteq \Gamma$
  kang winates bisa dicukupi.
\end{defn}

\begin{thm}[Teorema Kekompakan]
\ollabel{thm:compactness}
Pratelan-pratelan sabanjure bener kanggo saben himpunan ukara
$\Gamma$ lan ukara~$!A$. Cathetan koreksi sumber OLPL-053:
sumber Inggris kanthi kliru nyebut loro-lorone minangka ukara; ing
kene Gamma dibalekake minangka himpunan ukara, kaya kang dibutuhake
dening rong pratelan sabanjure.
\begin{enumerate}
  \item $\Gamma \Entails !A$ yen lan mung yen ana $\Gamma_0
    \subseteq \Gamma$ kang winates nganti $\Gamma_0 \Entails !A$.
  \item $\Gamma$ bisa dicukupi yen lan mung yen bisa dicukupi kanthi
    winates.
\end{enumerate}
\end{thm}

\begin{proof}
Kita mbuktekake (2). Yen $\Gamma$ bisa dicukupi, mula ana
\iftag{FOL}{!!a{structure}~$\Struct{M}$}{!!a{valuation}~$\pAssign{v}$}
nganti $\iftag{FOL}{\Sat{M}}{\pSat{v}}{!A}$ kanggo kabeh $!A \in
\Gamma$. Mesthi wae, $\iftag{FOL}{\Struct{M}}{\pAssign{v}}$ iki uga
nyukupi saben subhimpunan winates saka~$\Gamma$, mula $\Gamma$ bisa
dicukupi kanthi winates.

Saiki upamana $\Gamma$ bisa dicukupi kanthi winates. Mula saben
subhimpunan winates~$\Gamma_0 \subseteq \Gamma$ bisa dicukupi. Miturut
kesahihan
(\iftag{FOL}{%
  \tagrefs{prfAX/{fol:axd:sou:cor:consistency-soundness},
    prfSC/{fol:seq:sou:cor:consistency-soundness},
    prfND/{fol:ntd:sou:cor:consistency-soundness},
    prfTab/{fol:tab:sou:cor:consistency-soundness}}}{%
  \tagrefs{prfAX/{pl:axd:sou:cor:consistency-soundness},
    prfSC/{pl:seq:sou:cor:consistency-soundness},
    prfND/{pl:ntd:sou:cor:consistency-soundness},
    prfTab/{pl:tab:sou:cor:consistency-soundness}}}),
saben subhimpunan winates konsisten. Banjur $\Gamma$ dhewe kudu
konsisten miturut
\iftag{FOL}{%
  \tagrefs{prfAX/{fol:axd:ptn:prop:proves-compact},
    prfSC/{fol:seq:ptn:prop:proves-compact},
    prfND/{fol:ntd:ptn:prop:proves-compact},
    prfTab/{fol:tab:ptn:prop:proves-compact}}}{%
  \tagrefs{prfAX/{pl:axd:ptn:prop:proves-compact},
    prfSC/{pl:seq:ptn:prop:proves-compact},
    prfND/{pl:ntd:ptn:prop:proves-compact},
    prfTab/{pl:tab:ptn:prop:proves-compact}}}.
Miturut kelengkapan (\olref[cth]{thm:completeness}), amarga
$\Gamma$~konsisten, himpunan kasebut bisa dicukupi.
\end{proof}

\tagprob{FOL}
\begin{prob}
Buktekna (1) saka \olref[fol][com][com]{thm:compactness}.
\end{prob}
\tagendprob

\tagprob{notFOL}
\begin{prob}
Buktekna (1) saka \olref[pl][com][com]{thm:compactness}.
\end{prob}
\tagendprob

\iftag{FOL}{%
\begin{ex}
Ing saben model~$\Struct{M}$ saka sawijining teori~$\Gamma$, saben
term katutup~$t$ mesthi nuduhake !!a{element} saka~$\Domain{M}$. Apa kita bisa
njamin yen saben !!{element} saka~$\Domain{M}$ uga dituduhake dening
sawijining term katutup? Kanthi tembung liya, apa ana teori~$\Gamma$ kang
kabeh domaine dicakup dening term katutup? Teorema kekompakan nuduhake yen jawabane ora yen
$\Gamma$ nduweni model tanpa wates. Katrangane mangkene: Upamana
$\Struct{M}$ minangka model tanpa wates saka~$\Gamma$, lan upamana
$c$ minangka !!a{constant} kang ora ana ing basa~$\Gamma$. Upamana
$\Delta$ minangka himpunan kabeh ukara $\eq/[c][t]$ kanggo saben term katutup
$t$ ing basa~$\Lang{L}$ saka~$\Gamma$, yaiku (cathetan koreksi sumber OLPL-743: subskrip nol ing rumus sabanjure nuduhake term katutup supaya unsur Delta tetep ukara),
  \[
  \Delta = \Setabs{\eq/[c][t]}{t \in {\Trm[L]}_0}.
  \]
Sawijining subhimpunan winates saka $\Gamma \cup \Delta$ bisa ditulis
minangka $\Gamma' \cup \Delta'$, kanthi $\Gamma' \subseteq \Gamma$ lan
$\Delta' \subseteq \Delta$. Amarga $\Delta'$ winates, himpunan iki
mung bisa ngemot term kang cacahe winates. Jupuken $a \in \Domain{M}$
minangka !!a{element} saka $\Domain{M}$ kang ora dituduhake dening
siji wae saka term-term kasebut, lan upamana $\Struct{M'}$ minangka
!!{structure} kang padha karo $\Struct{M}$, nanging uga
$\Assign{c}{M'} = a$. Amarga $a \neq \Value{t}{M}$ kanggo kabeh~$t$
kang muncul ing~$\Delta'$, $\Sat{M'}{\Delta'}$. Amarga
$\Sat{M}{\Gamma}$, $\Gamma' \subseteq \Gamma$, lan $c$ ora muncul
ing~$\Gamma$, kita uga duwe $\Sat{M'}{\Gamma'}$. Dadi,
$\Sat{M'}{\Gamma' \cup \Delta'}$ kanggo saben subhimpunan winates
$\Gamma' \cup \Delta'$ saka $\Gamma \cup \Delta$. Mula saben
subhimpunan winates saka $\Gamma \cup \Delta$ bisa dicukupi. Miturut
kekompakan, $\Gamma \cup \Delta$ dhewe bisa dicukupi. Dadi,
$\Sat{M}{\Gamma \cup \Delta}$ kanggo sawijine model M. Saben
$\Struct{M}$ kaya mangkono
minangka model saka~$\Gamma$, nanging domaine ora dicakup dening term katutup
saka basa asline, amarga $\Value{c}{M} \neq \Value{t}{M}$ kanggo kabeh
term katutup~$t$ saka~$\Lang{L}$.
\end{ex}

\begin{ex}
Gatekna !!a{language} $\Lang{L}$ kang ngemot !!{predicate}~$<$,
!!{constant}s $\Obj{0}$, $\Obj{1}$, lan !!{function}s $+$, $\times$,
lan $-$. Upamana $\Gamma$ minangka himpunan kabeh !!{sentence}s ing
!!{language} iki kang bener ing !!{structure}~$\Struct{Q}$ kanthi
domain~$\Rat$ lan tafsiran kang lumrah. $\Gamma$~yaiku himpunan kabeh
!!{sentence}s saka~$\Lang{L}$ kang bener ngenani wilangan rasional.
Mesthi wae, ing~$\Rat$ (lan malah ing~$\Real$), ora ana wilangan~$r$
kang luwih gedhe tinimbang~$0$ nanging luwih cilik tinimbang $1/k$
kanggo kabeh $k \in \PosInt$. Yen ana, wilangan kaya mangkono bakal
dadi \emph{infinitesimal:} ora nol, nanging cilik tanpa wates.
Teorema kekompakan bisa digunakake kanggo nuduhake yen ana model
saka~$\Gamma$ kang ngemot infinitesimal. Ing basa kita ora ana
!!a{function} kanggo operasi mbagi (mbagi nganggo nol ora katetepake,
lan !!{function}s kudu ditafsirake dening fungsi total). Nanging, kita
isih bisa nyatakake yen $r < 1/k$, amarga iki bener yen lan mung yen
$r \cdot k < 1$. Saiki upamana $c$ minangka !!{constant} anyar lan
upamana $\Delta$ yaiku
\[
\{\Obj{0} < c\}
\cup \Setabs{ c \times \num k < \Obj{1} }{k \in \PosInt}
\]
(ing kene $\num{k} = (\Obj{1} + (\Obj{1} + \dots + (\Obj{1} +
\Obj{1})\dots))$ kanthi $k$~$\Obj{1}$). Kanggo saben subhimpunan
winates~$\Delta_0$ saka~$\Delta$, ana~$K$ positif nganti kabeh !!{sentence}s
$c \times \num{k} < \Obj{1}$ ing~$\Delta_0$ nduweni $k < K$. Yen
kita nggedhekake $\Struct{Q}$ dadi~$\Struct{Q'}$ kanthi
$\Assign{c}{Q'} = 1/K$, kita duwe
$\Sat{Q'}{\Gamma_0 \cup \Delta_0}$ kanggo saben $\Gamma_0 \subseteq
\Gamma$ kang winates, mula $\Gamma \cup \Delta$ bisa dicukupi kanthi
winates (Gladhen: buktekna iki kanthi rinci). Miturut kekompakan,
$\Gamma \cup \Delta$ bisa dicukupi. Saben model~$\Struct{S}$ saka
$\Gamma \cup \Delta$ ngemot infinitesimal, yaiku~$\Assign{c}{S}$.
\end{ex}
}{}

\tagprob{FOL}
\begin{prob}
Ing model standar aritmetika~$\Struct{N}$, ora ana !!{element}~$k \in
\Domain{N}$ kang nyukupi saben formula~$\num{n} < x$ (ing kene
$\num{n}$ yaiku $\Obj{0}^{\prime\dots\prime}$ kanthi $n$~$\prime$).
Gunakna teorema kekompakan kanggo nuduhake yen himpunan
!!{sentence}s ing basa aritmetika kang bener ing model standar
aritmetika $\Struct{N}$ uga bener ing !!a{structure}~$\Struct{N'}$
kang ngemot !!a{element} kang \emph{pancen} nyukupi saben formula
$\num{n} < x$.
\end{prob}
\tagendprob

\iftag{FOL}{
\begin{ex}
Kita ngerti yen logika sepisanan kanthi !!{identity} bisa nyatakake
yen ukuran !!{domain} kudu nduweni ukuran minimal tartamtu: Ukara
$!A_{\ge n}$ (kang nyatakake “ana paling ora $n$ objek kang beda”)
mung bener ing struktur kang $\Domain{M}$ ngemot paling ora~$n$ objek.
Dadi yen kita njupuk
  \[
  \Delta = \Setabs{!A_{\ge n}}{n \ge 1}
  \]
mula saben model saka $\Delta$ kudu tanpa wates. Mula, kita bisa
njamin yen sawijining teori mung nduweni model tanpa wates kanthi
nambahake~$\Delta$ marang teori kasebut: model saka $\Gamma \cup
\Delta$ yaiku kabeh lan mung model tanpa wates saka~$\Gamma$.

Dadi logika sepisanan bisa nyatakake sifat tanpa wates. Nanging,
teorema kekompakan nuduhake yen logika sepisanan ora bisa nyatakake
sifat winates. Upamana ana sawijine himpunan ukara $\Lambda$ kang bisa
dicukupi ing kabeh lan mung !!{structure}s winates. Banjur $\Delta
\cup \Lambda$ bisa dicukupi kanthi winates. Apa sebabe? Upamana
$\Delta' \cup \Lambda' \subseteq \Delta \cup \Lambda$ winates kanthi
$\Delta' \subseteq \Delta$ lan $\Lambda' \subseteq \Lambda$. Cathetan
koreksi sumber OLPL-054: yen subhimpunan Delta-prime kasebut kosong,
jupuken n minangka siji; yen ora kosong, gunakna angka paling gedhe
kaya ing ukara sabanjure. Ing perkara kang ora kosong, upamana $n$
minangka wilangan paling gedhe nganti $!A_{\ge n} \in \Delta'$; ing
perkara kosong, n wis dipilih minangka siji. Amarga $\Lambda$ bisa
dicukupi ing kabeh struktur winates, himpunan kasebut nduweni
model~$\Struct{M}$
kanthi !!{element}s kang cacahe winates nanging $\ge n$. Nanging banjur
$\Sat{M}{\Delta' \cup \Lambda'}$. Miturut kekompakan, $\Delta \cup
\Lambda$ nduweni model tanpa wates, kang mbantah panganggep yen
$\Lambda$ mung bisa dicukupi ing !!{structure}s winates.
\end{ex}
}{}

\end{document}
